\documentclass[11pt]{article}

\usepackage[utf8]{inputenc}
\usepackage[T1]{fontenc}
\usepackage{lmodern}
\usepackage[margin=1in]{geometry}
\usepackage{booktabs}
\usepackage{amsmath}
\usepackage{microtype}
\usepackage{graphicx}
\usepackage{subcaption}

\title{Natural Language Processing - Ex5\\\large Sentiment Analysis with Neural Models}
\author{Nurit Tolkowsky and Shaul Tolkowsky}
\date{}

\newcommand{\modelfig}[2]{%
\begin{figure}[h]
\centering
\begin{subfigure}{0.49\textwidth}\centering
\includegraphics[width=\textwidth]{#1_loss.png}\caption{Train/validation loss}
\end{subfigure}\hfill
\begin{subfigure}{0.49\textwidth}\centering
\includegraphics[width=\textwidth]{#1_acc.png}\caption{Train/validation accuracy}
\end{subfigure}
\caption{#2}
\end{figure}}

\begin{document}
\maketitle

We trained all four models on the Stanford Sentiment Treebank with the
hyperparameters required by the exercise. The three full-data models
(Sections 6--8) were trained on the full sentences together with their
sub-phrases (73{,}570 training items); the fine-tuned Transformer (Section 9)
was trained on 2{,}500 randomly sampled full sentences only. For each model we
report the train/validation loss and accuracy curves, the final test loss and
accuracy, and the accuracy on the two special test subsets (62 negated-polarity
sentences and 50 rare-word sentences).

\section*{Section 6 - Simple Log-Linear Model (one-hot)}

\modelfig{log_linear_onehot}{Section 6 - one-hot log-linear.}

\begin{center}
\begin{tabular}{lc}
\toprule
Quantity & Value \\
\midrule
Test loss & 0.4237 \\
Test accuracy & 0.8233 \\
Negated-polarity accuracy (62 ex.) & 0.6452 \\
Rare-words accuracy (50 ex.) & 0.5400 \\
\bottomrule
\end{tabular}
\end{center}

\section*{Section 7 - Transformer-Embedding Log-Linear Model}

\modelfig{log_linear_transformer}{Section 7 - Transformer-embedding log-linear.}

\begin{center}
\begin{tabular}{lc}
\toprule
Quantity & Value \\
\midrule
Test loss & 0.3967 \\
Test accuracy & 0.8368 \\
Negated-polarity accuracy (62 ex.) & 0.6290 \\
Rare-words accuracy (50 ex.) & 0.8400 \\
\bottomrule
\end{tabular}
\end{center}

\section*{Section 8 - Bi-LSTM Model}

\modelfig{lstm_transformer}{Section 8 - bi-LSTM with frozen Transformer embeddings.}

\begin{center}
\begin{tabular}{lc}
\toprule
Quantity & Value \\
\midrule
Test loss & 0.4266 \\
Test accuracy & 0.8514 \\
Negated-polarity accuracy (62 ex.) & 0.6129 \\
Rare-words accuracy (50 ex.) & 0.7800 \\
\bottomrule
\end{tabular}
\end{center}

\section*{Section 9 - Fine-Tuned Transformer Model (distilroberta)}

\modelfig{transformer_finetuned}{Section 9 - fine-tuned distilroberta.}

\begin{center}
\begin{tabular}{lc}
\toprule
Quantity & Value \\
\midrule
Test loss & 0.2756 \\
Test accuracy & 0.8919 \\
Negated-polarity accuracy (62 ex.) & 0.7419 \\
Rare-words accuracy (50 ex.) & 0.7400 \\
\bottomrule
\end{tabular}
\end{center}

\section{Section 10 - Comparing the Models}

\subsection*{10.1 \quad One-hot vs.\ Transformer-embedding log-linear}

The two log-linear models share the same architecture and differ only in the
input representation, so it is not surprising that their overall accuracy is
close (test 0.8233 vs.\ 0.8368, validation essentially tied). Both average the
word vectors and reduce the sentence to a bag of words, which caps how well
either can do.

The clear difference is on the rare-words subset (0.5400 vs.\ 0.8400). This
makes sense: the one-hot model gives every rare word its own coordinate that
was barely updated during training, so it carries almost no signal. In the
Transformer embedding space, on the other hand, a rare word most likely lands
in the same region as more common words of similar meaning, so the sentiment
signal learned for those common words generalizes to the rare one. The
embedding effectively turns "rare" into "close to something familiar," and that
is exactly the kind of generalization the one-hot representation cannot
provide.

\subsection*{10.2 \quad Adding the LSTM and the fine-tuned Transformer}

Both of the more expressive models beat both log-linear models. By test
accuracy the ordering is
\[
\text{fine-tuned (0.8919)} > \text{bi-LSTM (0.8514)} > \text{transf-avg (0.8368)} > \text{one-hot (0.8233)}.
\]
The bi-LSTM improves on the log-linear models because it reads the embeddings in
order (in both directions) instead of averaging them, so it can use word order
and composition.

Its training curves also show a point worth noting: the model overfits in the
\emph{loss} (validation loss rises after epoch 2) but not in the
\emph{accuracy} (validation accuracy stays roughly flat). The reason is that
accuracy only looks at which side of $0.5$ a prediction falls on, while the BCE
loss penalizes confidence. As training continues the model becomes more
certain, so its mistakes turn into confidently-wrong predictions and the loss
grows, even though the predictions do not actually flip sides. In other words,
the model is only becoming more certain about the mistakes it was already
making.

The fine-tuned Transformer is the best model by a clear margin, and it reached
this on only 2{,}500 sentences - roughly 30 times less data than the 73{,}570
items the other three models saw, and without any sub-phrases. Despite that
handicap it still wins, which is the strength of large-scale pretraining: the
model already encodes strong general-purpose representations, so a small amount
of task data is enough to adapt it. If anything its 0.8919 is a lower bound,
since with the full training set we would expect it to do at least as well.

This overall ordering does not hold everywhere, though. On the rare-words
subset the fine-tuned Transformer drops below the Transformer-embedding
log-linear model (0.7400 vs.\ 0.8400) - even though both use the same pretrained
distilroberta embeddings and differ only in whether those embeddings are
fine-tuned or frozen. We discuss this reversal in Section 10.3.

\subsection*{10.3 \quad The two special subsets}

\paragraph{Negated polarity.}
The fine-tuned Transformer is clearly best (0.7419); the other three cluster
together and much lower (one-hot 0.6452, transf-avg 0.6290, bi-LSTM 0.6129).
Negation requires understanding that a word like "not" flips the polarity of
what follows. The two averaging models cannot represent this at all, since they
sum word vectors and lose the interaction between the negator and the word it
negates. Only the fine-tuned Transformer, whose self-attention lets every word
condition on the negator, handles negation noticeably better. The bi-LSTM is in
fact the lowest of the four here, but with only 62 examples the differences
among the bottom three are small and within noise.

\paragraph{Rare words.}
This is the one place the fine-tuned model does \emph{not} win:
\[
\text{transf-avg (0.8400)} > \text{bi-LSTM (0.7800)} > \text{fine-tuned (0.7400)} \gg \text{one-hot (0.5400)}.
\]
The one-hot model is worst by a wide margin for the reason given in Section
10.1. The interesting part is that fine-tuning actually \emph{hurts} rare-word
accuracy relative to the frozen embedding. We read this as fine-tuning
destroying exactly the property that made rare words work in Section 10.1. In
the frozen embedding a rare word sits next to familiar words of similar
meaning, so its sentiment signal is generalizable; but when we fine-tune the
whole model on only 2{,}500 sentences, the weights move to fit that small
distribution and the rare word's neighbourhood structure is disrupted - the
rare word no longer benefits from sharing latent space with the common words.

This is a known danger of fine-tuning an entire pretrained model (it tends to
forget the structure it came with), and it is why for tasks like this one it is
often preferable to adapt the model more lightly: adding a LoRA adapter,
fine-tuning only the last layer, or unfreezing only a randomly chosen subset of
the weights. These keep most of the pretrained representation intact - including
the rare-word neighbourhood structure - while still adapting the model to the
task. (With only 50 examples in this subset some of the gap is also noise, but
the direction is consistent with this explanation.)

\section{Section 11 - Counting Parameters}

\paragraph{Setup / notation.}
\begin{itemize}
\item $d$ = model dimension (width of a single token's vector).
\item $V$ = vocabulary size.
\item $n$ = context length (how many tokens are fed in at once).
\item $L$ = number of layers.
\item A \emph{parameter} is one individual learned number (one entry of a
weight matrix). The parameter count is how many numbers must be stored and
learned.
\end{itemize}

\subsection*{11.1 \quad Transformer}

\paragraph{(a) Token embedding matrix.}
The embedding is a $V \times d$ lookup table (one $d$-dimensional vector per
vocabulary token), so it has
\[
V \cdot d \quad\text{parameters.}
\]

\paragraph{(b) Whole encoder (embedding $+$ $L$ attention layers).}
Each attention layer has four projections $W_Q, W_K, W_V, W_O$, each $d \times
d$, giving $4d^2$ per layer and $4Ld^2$ across $L$ layers. Adding the embedding,
\[
\boxed{\,V \cdot d + 4 L d^2\,}.
\]

\paragraph{Which part dominates? ($V = 30{,}000$, $d = 512$).}
\begin{itemize}
\item Embedding: $V d = 30{,}000 \times 512 \approx 15.4\text{M}$.
\item Attention: $4 d^2 = 4 \times 512^2 \approx 1.05\text{M}$ per layer.
\end{itemize}
The comparison reduces to $V$ vs.\ $4Ld$: that is $30{,}000$ vs.\ $\approx
2{,}048 \cdot L$, which are equal only around $L \approx 15$. For a typical small
$L$ the \textbf{embedding dominates}. Intuitively, the embedding carries a
factor of $V$ ($30{,}000$) in front of $d$, while each attention layer carries
only $4d \approx 2{,}000$; the huge vocabulary outweighs the $d^2$ of attention
unless many layers are stacked.

\paragraph{Does any part depend on $n$?}
\textbf{No} - neither $V d$ nor $4Ld^2$ contains $n$. The weights are
\emph{shared across all positions}: the embedding is one fixed $V \times d$
dictionary (a longer sequence just means more lookups in the same table, not
more rows), and the four $d \times d$ projections are applied to every token
with the same weights (a longer sequence reuses them more times, it does not
create new ones). So $n$ governs how many times the weights are \emph{used}
(runtime compute, and the $n^2$ attention map), not how many weights exist.

\subsection*{11.2 \quad RNN}

Replacing the $L$ attention layers with $L$ recurrent layers, where one
recurrent layer's time-shared transformation has a fixed parameter count $R$
(given as a black box), and keeping the embedding matrix (tokens still have to
become vectors):
\[
\boxed{\,V \cdot d + L R\,}.
\]
Here $V d$ is the same embedding matrix as the Transformer, and $LR$ is the
recurrent part ($R$ per layer $\times$ $L$ layers).

\paragraph{Does it depend on $n$?}
\textbf{No.} $R$ is fixed per layer regardless of sequence length: the recurrent
transformation is time-shared, so the same weights are applied at every
timestep and a longer sequence just means more timesteps reusing the same $R$
parameters. This is exactly what lets an RNN handle arbitrarily long sequences
with a fixed parameter budget.

\subsection*{11.3 \quad Characters as tokens}

Switching from word tokens to character tokens, $V$ shrinks a lot
($\sim$30{,}000 $\to$ $\sim$100) but $n$ grows ($\sim$5$\times$) to cover the
same amount of text.

\paragraph{(a) By what factor does the embedding matrix ($V \cdot d$) shrink?}
$d$ is unchanged, so only the $V$ ratio matters:
\[
30{,}000 / 100 = \textbf{\(\sim\)300\(\times\) smaller}.
\]
The $5\times$ growth in $n$ does not enter here, since $n$ is not in the
embedding term.

\paragraph{(b) Do the attention projections / recurrent parameters change?}
\textbf{No.} $4Ld^2$ depends only on $d$ and $L$, and $LR$ depends only on $L$
and $R$; none of these move when $V$ shrinks or $n$ grows. Since $n$ appears in
neither formula, growing it $5\times$ changes nothing.

\paragraph{Combining (a) and (b): does each model's total go up or down?}
\textbf{Down}, for both models. The embedding term - which \emph{dominated} the
total in Part 1 - collapses by $\sim$300$\times$, while everything else
(attention $4Ld^2$ / recurrent $LR$) stays flat, so the total drops
substantially.

In summary, growing $n$ by $5\times$ costs no parameters (only runtime compute),
while shrinking $V$ saves a large amount, so character-level tokenization makes
both models smaller in parameter count - the trade-off is paid in sequence
length (compute) rather than in weights.

\end{document}
